\exer{Geometrie -- EPAS}
\setcounter{exo}{0}
% Pole JC TD 4 Sup


On s'interesse a une Echelle Pivotante Automatique a commande Sequentielle (voir video sur site internet).
Ce systeme, concu et commercialise par la societe CAMIVA, est monte sur le chassis d'un camion de pompiers et permet de
deplacer une plate-forme, pouvant recevoir deux personnes et un brancard (charge maxi 270 kg), le plus rapidement
possible et en toute securite.

\begin{figure}[H]
\centering
\includegraphics[width=\linewidth]{994_01}
\end{figure}

\begin{figure}[H]
\centering
\includegraphics[width=\linewidth]{994_02}
\end{figure}



Le systeme est represente, dans la situation particuliere alpha = 0, sous forme de schema cinematique ci-contre.
Ce systeme est constitue de quatre solides, listes ci-dessous avec leur repere associe :
\begin{itemize}
\item chassis 0, R0 = (O, x0, y0, z0) ;
\item tourelle 1, R1 = (A, x1, y1, z1), avec y1 = y0 ;
\item berceau 2, R2 = (A, x2, y2, z2), avec z2 = z1 ;
\item echelle 3, R3 = (D, x3, y3, z3), avec B3 = B2 ;
\item plate-forme 6, R6 = (E, x6, y6, z6), avec z6 = z2.
\end{itemize}

On donne les parametres de mouvement :
alpha = (x0, x1), theta = (x1, x2), beta = (x2, x6), CD = lambda x2.

On donne les dimensions caracteristiques : OA = a x1 + b y1, AC = c x2, DE = e x6.

L'etude se fait pendant la phase de dressage. Pendant cette phase, la tourelle 1 est fixe par rapport au chassis 0 (alpha = cte) et le berceau 2 s'incline de plus en plus alors que l'echelle 3 se deploie. On suppose que le chassis 0 est ancre sur le sol
parfaitement horizontal grace aux bras de stabilisation.

On donne ci-dessous un extrait du cahier des charges :
\begin{itemize}
\item exigence 1 : pendant la phase de dressage, la plate-forme 6 doit rester en permanence horizontale afin d'assurer la
securite des personnes qui sont embarquees.
\item exigence 2 : lors d'une intervention, le point E a l'extremite de la plate-forme 6 doit se deplacer verticalement le long de la facade d'un immeuble (L = cte et y variable au cours du temps).
\end{itemize}

\begin{obj}
Determiner la loi de commande en position de chacune des chaines d'energie-puissance afin de
respecter les exigences du cahier des charges.
\end{obj}

\question{Sur le schema cinematique, repasser chaque solide d'une couleur differente.}
\ifprof
\begin{corrige}
\end{corrige}
\else
\fi

\question{Realiser le graphe des liaisons. Preciser le parametre de mouvement associe a chaque liaison. Preciser les
eventuelles egalites de base.}
\ifprof
\begin{corrige}
\end{corrige}
\else
\fi

\question{Realiser les figures de changement de base illustrant les parametres de mouvement angulaires.}
\ifprof
\begin{corrige}
\end{corrige}
\else
\fi

\question{Preciser en le justifiant, le mouvement de 6/1, afin de respecter l'exigence Ex1 du cahier des charges. Que peut-on dire de la base B6 = (x6, y6, z6) ? Donner la relation entre beta et gamma afin de garantir ce mouvement ?}
\ifprof
\begin{corrige}
\end{corrige}
\else
\fi

\question{Donner, en fonction des parametres de mouvements et des dimensions caracteristiques du mecanisme,
l'expression du vecteur position AE du point E dans le repere 1.}
\ifprof
\begin{corrige}
\end{corrige}
\else
\fi

\question{Donner, en fonction de L et y, une autre expression de ce vecteur lorsque l'exigence Ex2 est respectee.}
\ifprof
\begin{corrige}
\end{corrige}
\else
\fi

\question{En deduire deux relations mathematiques, traduisant l'exigence Ex2 du cahier des charges, qui lient les
parametres de mouvement, les dimensions caracteristiques, L et y.}
\ifprof
\begin{corrige}
\end{corrige}
\else
\fi

\question{Donner alors les lois de commande en position lambda = f(y) puis theta = f(y). Conclure.}
\ifprof
\begin{corrige}
\end{corrige}
\else
\fi

% BEGIN AUTO CORRIGE INCLUDE
\InclureCorrige{corrige.tex}
% END AUTO CORRIGE INCLUDE